\documentclass[11pt,a4paper]{article}

\usepackage[a4paper,margin=2.5cm]{geometry}
\usepackage{amsmath}
\usepackage{amssymb}
\usepackage{hyperref}

\title{\textbf{Topological Integrity Gravity (TIG):\\
A Spectral-Information-Theoretic Framework for Gravity and Quantum Fields}}

\author{Kai}

\date{\today}

\begin{document}

\maketitle

\begin{abstract}
We present Topological Integrity Gravity (TIG), a theoretical framework in which spacetime, quantum fields, and gravitational dynamics emerge from an underlying spectral structure. By extending the spectral action principle with an integrity functional, the theory defines a self-regulating system that suppresses ultraviolet divergences and avoids curvature singularities. In the low-energy limit, the model reproduces General Relativity with higher-order corrections. This approach provides a consistent connection between geometric and quantum descriptions without introducing ad hoc regularization schemes.
\end{abstract}

\section{Introduction}

General Relativity (GR) describes gravity as a geometric property of spacetime and has been experimentally confirmed across a wide range of scales. However, it exhibits fundamental limitations:

\begin{itemize}
\item the presence of curvature singularities,
\item the absence of an intrinsic ultraviolet (UV) structure,
\item the lack of a direct connection to quantum field theory (QFT).
\end{itemize}

This work proposes a framework in which geometry and quantum fields arise from a common spectral structure subject to a stability constraint.

\section{Spectral Framework}

Spacetime is described by a spectral triple:
\begin{equation}
(\mathcal{A}, \mathcal{H}, D),
\end{equation}
where $\mathcal{A}$ is an algebra of observables, $\mathcal{H}$ a Hilbert space, and $D$ a Dirac operator encoding geometric information.

\section{Dynamical Principle}

The dynamics is defined by the action:
\begin{equation}
S = \mathrm{Tr}\, f\left(\frac{D}{\Lambda}\right) + \lambda \mathrm{Tr}\,\log\left(1 + \frac{D^2}{\Lambda^2}\right),
\end{equation}
where $\Lambda$ is a characteristic energy scale and $\lambda$ a dimensionless coupling parameter.

The first term represents the spectral dynamics, while the second term acts as an integrity functional that enforces stability.

\section{Low-Energy Limit}

In the regime $D/\Lambda \ll 1$, the spectral action admits an expansion:
\begin{equation}
S_{\mathrm{eff}} \sim \int d^4x \sqrt{-g} \left( R + \alpha R^2 + \dots \right),
\end{equation}
which recovers the Einstein-Hilbert action with higher-order corrections.

\section{Emergence of Fields}

Fluctuations of the Dirac operator,
\begin{equation}
D = D_0 + \delta D,
\end{equation}
generate effective field degrees of freedom, including scalar modes, fermions, and gauge fields.

\section{Stability and Ultraviolet Behavior}

The logarithmic term in the action suppresses large eigenvalues of $D$ and introduces an effective ultraviolet cutoff. This ensures bounded curvature configurations and prevents divergent behavior.

\section{Conclusion}

Topological Integrity Gravity provides a unified framework in which gravitational and quantum phenomena emerge from a common spectral structure. The inclusion of an integrity functional leads to a self-regulating system that naturally avoids singularities and introduces a consistent ultraviolet behavior.

\end{document}
